In modern data analysis, some naturally-arising matrices are too large to compute and store efficiently. A canonical example comes from \emph{kernel methods} in machine learning. Given $n$ data points $x_1,\dots,x_n$ and a positive semidefinite kernel function $k(\cdot,\cdot)$, an associated kernel matrix $\bA \in \mathbb{R}^{n\times n}$ is defined entrywise by
\[
\bA_{ij} = k(x_i,x_j).
\]
Although we can compute any single entry $\bA_{ij}$ on demand, computing and storing all $n^2$ entries is prohibitively expensive when $n$, the size of the dataset, is large. Thus: \textit{how can we perform matrix computations when the input is available only through limited entry queries?} 

A large body of work addresses this problem for tasks such as low-rank approximation, matrix-vector and matrix-matrix multiplication, least-squares regression problems, and spectral approximation. 

\todo{revisit this exposition}

A common approach is to work with a small sketch of the matrix obtained by sampling a subset of rows or columns. One natural and widely studied sampling rule is \emph{length-squared sampling}, in which a row or column is sampled with probability proportional to its squared $\ell_2$ norm. This rule appears, for example, in the randomized Kaczmarz method for solving linear systems \cite{needell_stochastic_2015} and in additive-error algorithms for low-rank approximation \cite{frieze_fast_2004}. However, a direct implementation of length-squared sampling must have access to the relevant row or column norms. If these norms are not already available, then forming the sampling distribution requires a preprocessing pass over the entire input matrix. For dense $n\times n$ matrices, this takes $\Theta(n^2)$ time.


% The problem is compounded by downstream computations involving $\bA$: many linear algebra routines, including eigendecomposition and solving least-squares, typically require $O(n^3)$ arithmetic operations.

% Fortunately, the kernel matrix $\bA$ often exhibits low-rank structure, meaning that it can be approximated by a much smaller matrix without significant loss of accuracy. Substituting $\bA$ with an appropriate low-rank approximation can yield speedups of several orders of magnitude in practice compared to direct methods, especially in kernel ridge regression and kernel spectral clustering \cite{chen_randomly_2025}. Given the practical significance, \textit{how can we compute accurate low-rank approximations for large kernel matrices efficiently?}

\section{Our Results}
The main technical contribution is an \(O(n)\)-time algorithm for length-squared sampling from an \(n \times n\) positive semidefinite (psd) matrix. Given entry-wise access to a psd matrix \(\mathbf{A}\), we sample column \(j\) with probability
\[
p_j = \frac{\|\mathbf{A}_{\ast, j}\|_2^2}{\|\mathbf{A}\|_F^2}.
\]
Our algorithm relies on psd structure and does not apply to arbitrary matrices. 

We also show that the query complexity of our algorithm is optimal up to constant factors. In particular, we prove an information-theoretic lower bound showing that any algorithm for length-squared sampling needs to query \(\Omega(n)\) off-diagonal entries in the worst case.

\todo{make sure this lower bound claim is precise}

We next show how this sampling primitive can be used in several low-rank and spectral approximation problems for positive semidefinite matrices. We consider additive-error low-rank approximation, robust low-rank approximation, and eigenvalue approximation. In each application, the sampler replaces a length-squared column sampling step that would otherwise require computing all column norms.

Finally, we include numerical experiments comparing length-squared column sampling with other column-selection methods for low-rank approximation. The goal of these experiments is to identify regimes where length-squared sampling is a competitive practical choice.

\section{Outline}
\label{sec:introduction-outline}
\todo{revisit outline at the end}

The remainder of this thesis is organized as follows. Chapter 2 introduces the mathematical background and notation used throughout the thesis. We review low-rank approximation, positive semidefinite matrices, the Nystr\"om approximation, and several sampling primitives that appear in randomized matrix approximation.

Chapter 3 contains the main technical contribution: a fast algorithm for length-squared sampling from a positive semidefinite matrix. We present the rejection-sampling procedure, prove that it samples columns from the correct distribution in expected $O(n)$ time, and give a matching lower bound showing that $\Omega(n)$ entry queries are necessary in the worst case.

Chapter 4 develops consequences of the sampling primitive for matrix approximation. We first show how the sampler yields a simple estimator for the Frobenius norm. We then apply the sampler to additive-error low-rank approximation and robust low-rank approximation. Finally, we discuss an application to eigenvalue approximation, where the Frobenius norm estimate is used as a subroutine.

Chapter 5 presents numerical experiments. We compare length-squared sampling with other column-selection methods for low-rank approximation, including uniform sampling and leverage score sampling. We identify regimes where length-squared sampling is a competitive practical choice.

Chapter 6 \todo{fill in} concludes the thesis and discusses directions for future work.

